\begin{thebibliography}{99}
\addcontentsline{toc}{chapter}{参考文献}
\bibitem{yosida-zh} 吉田耕作（K\={o}saku Yosida）.
\textbf{泛函分析}，第六版，中译本.
主线目录依据本地文件《泛函分析 第6版 (吉田耕作) 中译本.pdf》。
目录位于 PDF 第 10--15 页，正文印刷第 1 页位于 PDF 第 16 页。
\bibitem{yosida-en} K\={o}saku Yosida.
\textbf{Functional Analysis}, Sixth Edition. Springer.
英文辅助依据本地文件《泛函分析 第6版 英文原版.pdf》。
目录位于 PDF 第 11--15 页，正文印刷第 1 页位于 PDF 第 16 页。
\bibitem{schwartz} L. 施瓦尔茨（Laurent Schwartz）著，姚家燕译.
\textbf{广义函数论}，法兰西数学精品译丛.
本稿依据工作目录中同名中文扫描版；具体查阅位置见附录 \ref{app:references}。
\bibitem{gelfand} I. M. 盖尔芳特，G. E. 希烙夫（Gelfand--Shilov）.
\textbf{广义函数 III：微分方程理论的若干问题}.
本稿依据工作目录中同名中文扫描版。
\bibitem{alexandroff-hopf} P. Alexandroff (Alexandrov), H. Hopf.
\textbf{Topologie}, Vol.~I. Springer, 1935.
吉田第零章拓扑参考；原书总书目英文版印刷第 469 页。
\bibitem{bourbaki-topology} N. Bourbaki.
\textbf{Topologie G\'en\'erale}.
Actualit\'es Scientifiques et Industrielles, nos.~858, 916, 1029, 1045, 1084.
Hermann, 1940--1942（按吉田英文版印刷第 470 页的分册书目记录）。
现代合订重印本另见 \href{https://link.springer.com/book/10.1007/978-3-540-33982-3}{\emph{Topologie g\'en\'erale, Chapitres 1 \`a 4}, Springer, 2007}；重印本与原书引文的版次须区别。
\bibitem{kelley-topology} J. L. Kelley.
\textbf{General Topology}. D.~Van Nostrand, 1955.
吉田第零章拓扑参考；原书总书目英文版印刷第 475 页。
后由 Springer 再版为 Graduate Texts in Mathematics, Vol.~27, 1975；
\href{https://link.springer.com/book/9780387901251}{出版社书目页}注明原版出版于 1955 年。
\bibitem{halmos-measure} P. R. Halmos.
\textbf{Measure Theory}. D.~Van Nostrand, 1950.
吉田第零章测度参考；原书总书目英文版印刷第 473 页。
后由 Springer 重印，见 \href{https://link.springer.com/book/10.1007/978-1-4684-9440-2}{出版社书目页}。
\bibitem{saks-integral} S. Saks.
\textbf{Theory of the Integral}. Warszawa, 1937.
按吉田英文版印刷第 481 页的书目记录；原书章末列为 Saks [1]。
\end{thebibliography}
\section*{原书参考书目整理（占位）}
\phantomsection
\addcontentsline{toc}{section}{原书参考书目整理（占位）}
中译本印刷第 491--514 页（PDF 第 506--529 页）；
英文版印刷第 469--486 页（PDF 第 484--501 页）。
第零章涉及的 Alexandroff--Hopf、Bourbaki、Kelley、Halmos、Saks 五条已核对并收入上表，其余原书参考书目后续按实际引用逐条整理。
